


\section{Related Work}\label{sec:related}
In this section, we review existing works on self-evolving LLM agents and external memory or skill harnesses for agent experience.

\paragraph{Evolving and self-improving LLM agents.}
A growing line of work explores how LLM agents self-improve through interaction and feedback. Beyond local reflection methods that iteratively revise outputs based on scalar or verbal feedback~\cite{shinn2023reflexion,madaan2023self}, experience-based approaches attempt to extract reusable knowledge offline. For instance, some works mine accumulated trajectory pools to synthesize natural language insights~\cite{zhao2024expel}, interaction rules~\cite{chen2024automanual}, or transferable directories~\cite{ni2026trace2skill}. Similarly, a parallel track focuses on passively evolving higher-level agent artifacts, such as prompts~\cite{suzgun2025dynamic,lou2026autoharness}, reasoning banks~\cite{ouyang2025reasoningbank}, or memory structures~\cite{agrawal2025gepa,lee2026meta,zhang2025agentic,wei2025evo,lin2026position}. To provide more operational guidance, recent methods have shifted toward skill-centered evolution. This includes building executable code skills~\cite{wang2023voyager}, optimizing skill libraries via policy learning~\cite{xia2026skillrl,zhang2026memskill,ouyang2026skillos}, or curating skill knowledge bases and representations~\cite{yang2026autoskill,qiu2026autorefine,jiang2026xskill,wang2026skillx}. While closely related, these existing frameworks largely treat learned artifacts as passive repositories for free-form retrieval. Critically, prior methods often fail to rigorously validate their effectiveness under realistic online constraints, nor do they systematically isolate the underlying drivers of agent improvement. Our work addresses this gap by compiling single-shot experiences into an active skill harness, providing a structural scaffold to systematically investigate how LLM agents learn on the fly.


% A growing line of work studies how LLM agents improve through interaction, feedback, and accumulated experience. ReAct establishes a reasoning-and-acting paradigm for tool and environment interaction~\cite{yao2022react}, while reflection-based methods use feedback to revise behavior, such as Reflexion storing verbal reflections from failures and Self-Refine iteratively critiquing and improving model outputs~\cite{yao2022react， shinn2023reflexion,madaan2023self}. Beyond local revision, experience-based methods extract reusable knowledge from trajectories: ExpeL summarizes past experiences into natural-language insights, AutoManual converts interaction traces into rules and manuals, and Trace2Skill analyzes trajectory pools to build transferable skill directories~\cite{zhao2024expel,chen2024automanual,ni2026trace2skill}. Other works evolve higher-level agent artifacts, including prompts, harnesses, reasoning knowledge, and reusable guidance, as in GEPA, Dynamic Cheatsheet, Meta-Harness, AutoHarness, ReasoningBank, ACE, Evo-Memory, and A-Evolve~\cite{agrawal2025gepa,suzgun2025dynamic,lee2026meta,lou2026autoharness,ouyang2025reasoningbank,zhang2025agentic,wei2025evo,lin2026position}. \xy{split the reference}\xy{This paragraph includes too many aspects of evolution. Maybe keep this paragraph focused on classic evolving/self-improving agents, e.g., EvolveR-style methods, and discuss skill-centered evolution separately in the next paragraph.} Skill-centered methods provide a more operational form of evolution: Voyager builds executable code skills for open-ended exploration~\cite{wang2023voyager}; SkillRL evolves skill libraries with policy learning; SkillOS studies skill curation with a frozen executor and learned curator~\cite{xia2026skillrl,ouyang2026skillos}; and AutoSkill, AutoRefine, XSkill, SkillX, and MemSkill explore editable skills, skill/subagent patterns, experience-skill separation, skill knowledge bases, and evolvable memory skills~\cite{yang2026autoskill,qiu2026autorefine,jiang2026xskill,wang2026skillx,zhang2026memskill}. These works show that agents can evolve through reflections, prompts, rules, skills, memories, and repositories. Our work studies a challenging online setting, where a frozen agent incrementally updates a reusable skill harness from sequential task batches as they happen. \yc{This paragraph covers (too) many related works (and so long). Since these works are close to the topic of this paper, could we go beyond the final sentence that says “Our work studies a challenging online setting.”? It would be helpful to more specifically state how this paper differs from them.}

\paragraph{Memory and skill harnesses for agent experience.}

Memory provides a fundamental external medium for preserving agent experience. Existing memory systems maintain long-term observation histories~\cite{zhong2024memorybank,yu2025memagent}, organize memories through adaptive structures~\cite{xu2025mem}, and explore methods for updating and retrieving evidence across modalities~\cite{huang2026rethinking,bei2026mem,liu2026omnimem}. However, these approaches largely treat memory as a passive information store. Skills differ fundamentally because they encode operational knowledge. They shift from merely recording what happened to specifying exactly how an agent should act through manuals~\cite{chen2024automanual}, verify operations via organized directories~\cite{ni2026trace2skill}, or recover from failures using skill repositories~\cite{ouyang2026skillos,wang2026skillx,zhang2026mmskills}. This procedural view elevates skills into an active harness. Rather than merely supplying historical context, a harness acts as an external scaffold that structurally shapes future planning, taking forms such as pattern repositories~\cite{ning2026code,qiu2026autorefine}, reasoning banks~\cite{ouyang2025reasoningbank}, and dynamic guidance files~\cite{suzgun2025dynamic,lee2026meta,liu2026adaptive,lin2026harness}. Building on this distinction, our work conceptualizes the skill harness not merely as a performance-enhancing artifact, but as an interpretable medium to systematically study continuous online adaptation. To navigate the high noise of real-world executions, \algname{} compiles single-shot reflections into a structured harness of general and topic skills. While this separation efficiently filters out task-specific artifacts, its primary purpose in our work is to provide a transparent scaffold. This explicitly allows us to systematically isolate and evaluate the core drivers, such as evolver designs and feedback types, that dictate how frozen LLM agents self-improve on the fly.


% Memory provides an external medium for preserving and reusing agent experience. MemoryBank maintains long-term memories for persistent interactions, MemAgent studies memory-augmented agent behavior, and A-Mem organizes memories through adaptive memory structures~\cite{zhong2024memorybank,yu2025memagent,xu2025mem}. Recent memory systems and surveys further study how agent memories are stored, retrieved, updated, scheduled, and organized across modalities~\cite{huang2026rethinking,bei2026mem,liu2026omnimem}. These works mainly treat memory as an information store, preserving observations, histories, preferences, or retrieved evidence. Skills differ from general memory because they encode procedural knowledge that can directly guide future actions. Manuals, skill directories, SkillRepos, and skill knowledge bases move from storing what happened toward specifying how an agent should act, verify, recover, or use tools in later tasks~\cite{chen2024automanual,ni2026trace2skill,ouyang2026skillos,wang2026skillx,zhang2026mmskills}. This procedural view connects memory and skills to agent harnesses: a harness is not only a memory store, but an external structure that shapes future task solving. Prior work explores several forms of such structures, including skill libraries, experience-pattern repositories, reasoning banks, dynamic guidance files, and adaptive harnesses~\cite{qiu2026autorefine,ouyang2025reasoningbank,suzgun2025dynamic,lee2026meta}. Our work treats the skill harness as the medium of online adaptation from diverse and realistic scenarios. The agent compiles noisy execution contexts from task-specific reflection to topic and general skills, where general skills capture transferable patterns and topic skills preserve localized operational knowledge for recurring task types or contexts.

